\documentclass[11pt]{article}
\usepackage[margin=1in]{geometry}
\usepackage{amsmath,amssymb,amsthm,mathtools}
\usepackage[hidelinks,pdfusetitle]{hyperref}
\usepackage{microtype}
\newcommand{\DOI}[1]{\mbox{\href{https://doi.org/#1}{\nolinkurl{#1}}}}
\newcommand*{\PAPERDOI}{10.5281/zenodo.24412193}
\newcommand*{\SUPPLEMENTDOI}{10.5281/zenodo.24412195}
\theoremstyle{plain}
\newtheorem{theorem}{Theorem}
\newtheorem{lemma}{Lemma}
\newtheorem{proposition}{Proposition}
\newtheorem{corollary}{Corollary}
\theoremstyle{definition}
\newtheorem{definition}{Definition}
\DeclareMathOperator{\Tr}{Tr}
\DeclareMathOperator{\rank}{rank}
\newcommand{\Nodd}{N_{\mathrm{odd}}}
\newcommand{\Dphys}{D_{\mathrm{phys}}}
\newcommand{\Daug}{D_{\mathrm{aug}}}
\newcommand{\Dop}{\Delta}
\newcommand{\dist}{\operatorname{dist}}
\title{A fourth-power strong-coupling ceiling for electronic pair clustering
on hypercubic lattices in three and more dimensions}
\author{Davis Rens\\SuperMatics Inc.\\\href{https://orcid.org/0000-0001-5839-5197}{ORCID 0000-0001-5839-5197}}
\date{}
\begin{document}
\maketitle

\begin{abstract}
We prove exponential clustering of equal-time on-site electronic pair correlations
for the spinful Holstein model on the hypercubic torus $(\mathbb{Z}/M\mathbb{Z})^d$
in every dimension $d\ge2$, with explicit constants that track the dimension.
With $U=F^2/K$ the only hypothesis is
$\beta t^2/U\le 1/C(d)$, where $C(d)=442368\,d^2(d+1)^2$; no restriction on the
ratio $t/U$ is imposed, and the bound is uniform in the phonon frequency, the torus
side and the chemical potential. At distinct sites
\[
 \bigl|\langle\Dop_x^\dagger\Dop_y\rangle\bigr|\le A(d)\,q(d)^{|x-y|_1},
 \qquad A(d)=\frac{360(d+1)}{14d+29},\quad q(d)=\frac{5(2d-1)}{24(d+1)},
\]
with $\Dop_x=c_{x\downarrow}c_{x\uparrow}$. Hence a temperature below which pair
correlations need not cluster obeys $k_BT_{\mathrm{pair}}(d)\le C(d)\,t^2/U$.
In dimensions two and three the constants are explicit: $C(2)=15\,925\,248<2^{27}$,
which improves the released square-lattice constant $2^{27}$ by a factor $8.43$,
and $C(3)=63\,700\,992<2^{26}$, the first bound of this kind in dimension three.
The mechanism has two independent parts. The hopping majorant is controlled edge by
edge by an exact sum of squares on the sixteen-state bond Fock space, valid in every
particle-number sector and for arbitrarily large hopping, so no smallness condition
on $t/U$ arises. The combinatorial cost of connected bond supports is then paid by
injecting site animals into rooted subtrees of a Bethe lattice, which replaces
doubled-edge Euler tours and lowers the ceiling from order $d^5$ to order $d^4$.
The bounds concern a clustering temperature of a two-point function, not a critical
temperature, and no superconducting transition is claimed.
\end{abstract}

\section{Model and theorem}\label{sec:model}

\subsection{The spinful Holstein model on a hypercubic torus}

Fix an integer dimension $d\ge2$ and an integer side $M\ge4$. Let
$\Lambda=(\mathbb{Z}/M\mathbb{Z})^d$ be the hypercubic discrete torus of volume
$V=|\Lambda|=M^d$, with undirected nearest-neighbour edge set
\[
 E(\Lambda)=\bigl\{\{x,y\}\subset\Lambda:\ |x-y|_1=1\bigr\},
 \qquad |E(\Lambda)|=dV,
\]
where $|x-y|_1=\sum_{j=1}^d\min(|x_j-y_j|,M-|x_j-y_j|)$ is the torus graph
distance. Every site has physical coordination number $\Dphys(d)=2d$.

On the tensor product of the electronic Fock space with one oscillator per site, let
$c_{x\sigma},c_{x\sigma}^\dagger$ be canonical fermion operators with
$\sigma\in\{\uparrow,\downarrow\}$, and set
$n_{x\sigma}=c_{x\sigma}^\dagger c_{x\sigma}$,
$n_x=n_{x\uparrow}+n_{x\downarrow}\in\{0,1,2\}$. With hopping amplitude $t\ge0$,
oscillator coordinate $X_x$ and momentum $P_x$, mass $m>0$, spring constant $K>0$,
coupling $F\ne0$ and chemical potential $\mu\in\mathbb{R}$, the grand-canonical
Hamiltonian is
\begin{equation}\label{eq:hamiltonian}
 H_\Lambda(\mu)=-t\sum_{\{x,y\}\in E(\Lambda)}B_{\{x,y\}}
 +\sum_{x\in\Lambda}\left[\frac{P_x^2}{2m}+\frac{KX_x^2}{2}-FX_x(n_x-1)\right]
 -\mu\sum_{x\in\Lambda}(n_x-1),
\end{equation}
where
$B_{\{x,y\}}=\sum_{\sigma}(c_{x\sigma}^\dagger c_{y\sigma}+c_{y\sigma}^\dagger c_{x\sigma})$.
We write
\begin{equation}\label{eq:scales}
 U=\frac{F^2}{K}>0,\qquad \Omega=\hbar\sqrt{K/m}>0,\qquad
 \tau=k_BT>0,\qquad \beta=\frac{1}{\tau},
\end{equation}
all in units of energy. Each local oscillator has been centred, so completing the
square produces the uniform shift $-\frac U2\sum_x(n_x-1)^2$ and the atomic limit at
$\mu=0$ is symmetric about half filling. The chemical term is written as
$-\mu\sum_x(n_x-1)$; writing it as $-\mu\sum_x n_x$ instead, as in the released
two-dimensional record \cite{rens2d}, changes the grand canonical partition function
by the constant factor $e^{\beta\mu V}$ and leaves every correlation unchanged. There
is no direct electron-electron interaction.

The equal-time on-site singlet pair field is
\[
 \Dop_x=c_{x\downarrow}c_{x\uparrow},\qquad
 \Dop_x^\dagger=c_{x\uparrow}^\dagger c_{x\downarrow}^\dagger,\qquad
 \|\Dop_x\|_{\mathrm{op}}=1 .
\]
The field carries two electronic charges of the same sign and is therefore
conjugated by a squared displacement under the polaron transformation of
Section~\ref{sec:langfirsov}.

\begin{definition}[Clustering temperature]\label{def:tpair}
For $d\ge2$, let $T_{\mathrm{pair}}(d)$ be the infimum of the temperatures $\tau$
with the following property. For every $\tau'>\tau$, every $M\ge4$, every real
$\mu$ and every $\Omega>0$, there are constants $A_0,\xi>0$, independent of the
volume, of $\mu$ and of $\Omega$, such that
$|\langle\Dop_x^\dagger\Dop_y\rangle_{\beta',\mu,\Lambda}|
 \le A_0e^{-|x-y|_1/\xi}$ for all distinct $x,y\in\Lambda$, and such that the rows
of the correlation matrix are summable uniformly in $M$:
$\sup_{M\ge4}\max_{x\in\Lambda}\sum_{y\in\Lambda}
 |\langle\Dop_x^\dagger\Dop_y\rangle|<\infty$.
\end{definition}

$T_{\mathrm{pair}}$ is a clustering temperature. It is an upper threshold for
exponential decay of one two-point function, and it says nothing about the existence
of a phase transition at a critical temperature. Theorem~\ref{thm:main} below bounds
$T_{\mathrm{pair}}$ from above; no statement about a superconducting state, about a
transition temperature $T_c$, or about any material is made anywhere in this paper.

\subsection{Main theorem}

\begin{theorem}[Dimension-dependent pair-clustering ceiling]\label{thm:main}
Let $d\ge2$ and let $\Lambda=(\mathbb{Z}/M\mathbb{Z})^d$ be a hypercubic torus of
side $M\ge4$. Define
\begin{equation}\label{eq:constants}
 \begin{aligned}
 \Dphys(d)&=2d, & \Daug(d)&=2d+2, & L(d)&=36\,\Daug(d)=72(d+1),\\
 \varepsilon_0(d)&=+\infty, &
 C(d)&=\frac{256}{3}\,d^2L(d)^2=442368\,d^2(d+1)^2, & &
 \end{aligned}
\end{equation}
\begin{equation}\label{eq:constants2}
\begin{gathered}
 q(d)=\frac{5(2d-1)}{24(d+1)}<\frac{5}{12},\qquad
 A(d)=\frac{15}{1-q(d)}=\frac{360(d+1)}{14d+29}<\frac{180}{7},\\
 q_{\mathrm{sharp}}(d)=\frac{6\,903\,125\,(2d-1)}{40\,310\,784\,(d+1)}<q(d).
\end{gathered}
\end{equation}
Then for every $t\ge0$, every $U>0$, every $\Omega>0$, every $M\ge4$ and every
$\mu\in\mathbb{R}$, if
\begin{equation}\label{eq:hypothesis}
 \frac{\beta t^2}{U}\le\frac{1}{C(d)}
 \qquad\Longleftrightarrow\qquad
 k_BT\ge C(d)\,\frac{t^2}{U},
\end{equation}
then, for all distinct $x,y\in\Lambda$,
\begin{equation}\label{eq:mainbound}
 \bigl|\langle\Dop_x^\dagger\Dop_y\rangle_{\beta,\mu,\Lambda}\bigr|
 \le A(d)\,q_{\mathrm{sharp}}(d)^{|x-y|_1}
 \le A(d)\,q(d)^{|x-y|_1}.
\end{equation}
Moreover:
\begin{enumerate}
\item the rows are uniformly summable, with
\begin{equation}\label{eq:rowsum}
 \sup_{M\ge4}\max_{x\in\Lambda}\sum_{y\in\Lambda}
 \bigl|\langle\Dop_x^\dagger\Dop_y\rangle_{\beta,\mu,\Lambda}\bigr|
 \le 1+A(d)\left[\left(\frac{1+q(d)}{1-q(d)}\right)^d-1\right]<\infty;
\end{equation}
\item by the number-conserving pair-source construction of Section~\ref{sec:kp} and
the pressure transfer of Section~\ref{sec:canonical}, the bounds \eqref{eq:mainbound} and \eqref{eq:rowsum}
hold with the same constants under hypothesis \eqref{eq:hypothesis} in every
translation-invariant thermodynamic canonical limit at interior total density
$\rho\in(0,2)$, including balanced separately fixed spin densities
$\rho_\uparrow=\rho_\downarrow=\rho/2\in(0,1)$;
\item consequently, uniformly in $\Omega>0$, in the volume and in the chemical
potential, the clustering temperature of Definition~\ref{def:tpair} satisfies
\begin{equation}\label{eq:tpc}
 k_BT_{\mathrm{pair}}(d)\le C(d)\,\frac{t^2}{U}=442368\,d^2(d+1)^2\,\frac{t^2}{U}.
\end{equation}
\end{enumerate}
No condition on the ratio $t/U$ is imposed: hypothesis \eqref{eq:hypothesis} alone is
used, and the bound holds for arbitrarily strong coupling.
\end{theorem}

The explicit constants are collected in Table~\ref{tab:constants}. The values of
$C(d)$, $q(d)$, $q_{\mathrm{sharp}}(d)$ and $A(d)$, and the inclusion
$q_{\mathrm{sharp}}(d)<q(d)$, are certified by exact rational arithmetic for
$d=2,3,4,5$ (Section~\ref{sec:supplement}).

\begin{table}[t]
\centering
\small
\setlength{\tabcolsep}{5pt}
\begin{tabular}{@{}cccccccc@{}}
\hline
$d$ & $\Dphys=2d$ & $\Daug=2d+2$ & $L(d)$ & $C(d)$ & $\log_2C(d)$ & $q(d)$ & $A(d)$\\
\hline
$2$ & $4$ & $6$  & $216$ & $15\,925\,248$  & $23.925$ & $5/24$ & $360/19$\\
$3$ & $6$ & $8$  & $288$ & $63\,700\,992$  & $25.925$ & $25/96$ & $1440/71$\\
$4$ & $8$ & $10$ & $360$ & $176\,947\,200$ & $27.399$ & $7/24$ & $360/17$\\
$5$ & $10$& $12$ & $432$ & $398\,131\,200$ & $28.569$ & $5/16$ & $240/11$\\
\hline
\end{tabular}
\caption{Constants of Theorem~\ref{thm:main}. In every dimension the smallness
cutoff on $t/U$ is absent, $\varepsilon_0(d)=+\infty$. The sharp decay rate is
$q_{\mathrm{sharp}}(d)=0.17125$ for $d=2$, $0.21406$ for $d=3$, $0.23975$ for $d=4$
and $0.25687$ for $d=5$.}
\label{tab:constants}
\end{table}

Two comparisons fix the scale. In dimension two the released record
\cite{rens2d} proves the ceiling $k_BT_{\mathrm{pair}}\le 2^{27}t^2/U$ with decay
rate $5/64$ and prefactor $240/59$. Theorem~\ref{thm:main} gives
$C(2)=15\,925\,248<2^{27}=134\,217\,728$, a factor $8.43$ smaller, at the price of a
slower decay rate, $q(2)=5/24$ against $5/64$, and of a larger prefactor,
$A(2)=360/19$ against $240/59$. In dimension three,
\begin{equation}\label{eq:c3}
 C(3)=442368\cdot9\cdot16=63\,700\,992<2^{26}=67\,108\,864<2^{27},
\end{equation}
which is below the two-dimensional constant of the released result and, as far as we
know, is the first explicit pair-clustering ceiling for the three-dimensional
Holstein model that is stated in terms of the microscopic parameters alone.

\subsection{Baseline routes and the failure of a naive cubic extension}\label{sec:baselines}

Two technical ingredients enter Theorem~\ref{thm:main}: the edge-local sum of
squares of Section~\ref{sec:forms} in place of a conditional global form bound, and
Bethe-tree vertex-animal counting in place of doubled-edge Euler-tour counting. To
separate their effects we also record three baseline extensions, each a verbatim
transcription of a released square-lattice proof with the degrees promoted to
$2d$ and $2d+2$. Their proofs are in Section~\ref{sec:baselineproofs}.

\begin{proposition}[Edge-local squares with Euler-tour counting]\label{prop:a0}
Set $L_{\mathrm{ET}}(d)=24\,\Daug(d)^2=96(d+1)^2$,
$\varepsilon_0(d)=+\infty$ and
\begin{equation}\label{eq:a0constants}
 C_{\mathrm{ET}}(d)=128\,d\,L_{\mathrm{ET}}(d)^2
 =1\,179\,648\,d(d+1)^4=\Theta(d^5).
\end{equation}
If $\beta t^2/U\le1/C_{\mathrm{ET}}(d)$, then for all distinct $x,y$,
\begin{equation}\label{eq:a0bound}
 \bigl|\langle\Dop_x^\dagger\Dop_y\rangle\bigr|
 \le\frac{76}{5}\,q_{\mathrm{ET}}(d)^{|x-y|_1},
 \qquad q_{\mathrm{ET}}(d)=\frac{5d^2}{24(d+1)^2}<\frac{5}{24}.
\end{equation}
\end{proposition}

\begin{proposition}[Conditional global form with Euler-tour counting]\label{prop:b1}
Set $L_{\mathrm{CF}}(d)=2^{10}\Daug(d)^2=4096(d+1)^2$,
\begin{equation}\label{eq:b1constants}
\begin{gathered}
 \varepsilon_0(d)=\frac{1}{32\,d\,L_{\mathrm{CF}}(d)}=\frac{1}{2^{17}d(d+1)^2}
 =\Theta(d^{-3}),\\
 C_{\mathrm{CF}}(d)=128\,d\,L_{\mathrm{CF}}(d)^2=2^{31}d(d+1)^4=\Theta(d^5).
\end{gathered}
\end{equation}
If $t/U\le\varepsilon_0(d)$ and $\beta t^2/U\le1/C_{\mathrm{CF}}(d)$, then for all
distinct $x,y$,
\begin{equation}\label{eq:b1bound}
 \bigl|\langle\Dop_x^\dagger\Dop_y\rangle\bigr|
 \le\frac{9}{1-q_{\mathrm{CF}}(d)}\,q_{\mathrm{CF}}(d)^{|x-y|_1},
 \qquad q_{\mathrm{CF}}(d)=\frac{27d^2}{1024(d+1)^2}.
\end{equation}
\end{proposition}

\begin{proposition}[Conditional global form with Bethe-tree animals]\label{prop:b2}
Retain the conditional bound $t/U\le1/(16\Dphys(d))=1/(32d)$ with gap share
$\frac U4\Nodd(S)$ and the coarse fugacity bound $z_b/z_a<e$, but count vertex
animals. Set $L_{\mathrm{VA}}(d)=120\,\Daug(d)=240(d+1)$,
\begin{equation}\label{eq:b2constants}
\begin{gathered}
 \varepsilon_0(d)=\frac{1}{32\,d\,L_{\mathrm{VA}}(d)}=\frac{1}{7680\,d(d+1)}
 =\Theta(d^{-2}),\\
 C_{\mathrm{VA}}(d)=\frac{512}{3}\,d^2L_{\mathrm{VA}}(d)^2
 =9\,830\,400\,d^2(d+1)^2=\Theta(d^4).
\end{gathered}
\end{equation}
If $t/U\le\varepsilon_0(d)$ and $\beta t^2/U\le1/C_{\mathrm{VA}}(d)$, then pair
correlations decay exponentially at rate $q_{\mathrm{VA}}(d)<5/11$.
\end{proposition}

For $d=2,3,4,5$ the ceilings are ordered
$C_{\mathrm{ET}}(d)<C_{\mathrm{VA}}(d)<C_{\mathrm{CF}}(d)$, and each is certified
exactly. So vertex-animal counting improves only the exponent of $d$, by one power,
and at $d=3$ the unconditional edge-local estimate is then worth a further factor of
$200/9$ in the ceiling, from $C_{\mathrm{VA}}(3)=1\,415\,577\,600$ down to
$C(3)=63\,700\,992$.
Proposition~\ref{prop:c} explains why the Euler-tour counting cannot be pushed to
$2^{27}$ in dimension three.

\begin{proposition}[Cubic obstruction for Euler-tour counting]\label{prop:c}
In the canonical extension of the Euler-tour route at $d=3$, the Koteck\'y-Preiss
condition reads
\begin{equation}\label{eq:c19}
 \frac{15}{4}\,\frac{q_{\mathrm{src}}}{1-q_{\mathrm{src}}}\le1
 \qquad\Longleftrightarrow\qquad
 q_{\mathrm{src}}\le\frac{4}{19}.
\end{equation}
At $\beta t^2/U=2^{-27}$, for every choice of the scale $L_7>0$ the source activity
of that route satisfies the exact rational lower bound
\begin{equation}\label{eq:c675}
 q_{\mathrm{src}}(L_7)^2=\left(\frac{240}{L_7}\,
 e^{48L_7^2/2^{27}}\right)^2>\frac{675}{8192}
 >\frac{16}{361}=\left(\frac{4}{19}\right)^2 .
\end{equation}
Consequently no value of the scale admits the constant $2^{27}$ in that route in
dimension three, and its least admissible ceiling there exceeds $2^{28.33}$.
\end{proposition}

\section{Context and scope}\label{sec:context}

The Holstein model is the standard local electron-phonon model \cite{holstein}, and
the polaron transformation used below is the Lang-Firsov transformation
\cite{langfirsov}. The released square-lattice record \cite{rens2d}, with its
supplement \cite{rens2dsupp}, proves the two-dimensional ceiling
$k_BT_{\mathrm{pair}}\le2^{27}t^2/U$ together with the auxiliary estimates that the
present paper promotes to general dimension; here those estimates are used as
proved, with every dimension-dependent constant tracked.

Cluster expansions for quantum lattice systems at intermediate temperature were
developed by Fr\"ohlich, Rey-Bellet and Ueltschi \cite{fruhlich}, on the basis of the
polymer methods of Datta, Fern\'andez and Fr\"ohlich \cite{dff}, Borgs, Koteck\'y and
Ueltschi \cite{bku} and Borgs and Koteck\'y \cite{bkfermion}; the general criterion
for abstract polymer models is that of Koteck\'y and Preiss \cite{kp}, and
Ueltschi's Theorem 3.1 gives uniqueness and exponential clustering for the repulsive
spinful Hubbard model in its low-temperature domain, which includes small
$\beta t^2/\Delta$ and small $t/\Delta$ \cite{uelschiHubbard}. Ueltschi also develops
correlation-function bounds for cluster expansions \cite{uelschiCorrelations}. On a
bipartite lattice at particle-hole-symmetric half filling, a partial particle-hole
transformation on one spin maps the attractive Hubbard model to the half-filled
repulsive model and carries the on-site pair field to a staggered transverse spin.
Pair clustering at the $t^2/U$ scale is therefore already known for the attractive
Hubbard model at half filling. The high-frequency Holstein limit is the attractive
Hubbard model, but the two models differ away from that limit, and the present
statement is about the original oscillator degrees of freedom at every
$\Omega>0$. Charge-density-wave and strong-coupling expansion results for related
models are found in \cite{macris,freericks}.

The distinction claimed here is narrow: an explicit exponential bound on the
on-site electronic pair correlation of the original spinful Holstein model, on a
hypercubic torus of arbitrary dimension $d\ge2$, with constants that are uniform in
the phonon frequency and that are given in closed form in $d$, together with the
first such bound in dimension three below the released two-dimensional constant.
This is a statement about a two-point function of the model. It does not identify a
transition, does not exclude other forms of order, and makes no material prediction.

\section{The oscillator contraction is a contraction}\label{sec:langfirsov}

Apply the on-site Lang-Firsov unitary
\[
 \mathcal U=\exp\!\left(-\alpha_{\mathrm{LF}}\sum_{x\in\Lambda}
 (n_x-1)(b_x^\dagger-b_x)\right),
 \qquad
 \alpha_{\mathrm{LF}}=\frac{F}{\sqrt{2m\hbar\Omega^3}}
 =\sqrt{\frac{U}{2\hbar\Omega}},
\]
with $X_x=\sqrt{\hbar/(2m\Omega)}\,(b_x+b_x^\dagger)$. The transformed Hamiltonian
is $H_{\mathrm{LF}}=H_0+V_{\mathrm{LF}}$, where the atomic part is diagonal,
\[
 H_0=-\frac U2\sum_{x\in\Lambda}(n_x-1)^2-\mu\sum_{x\in\Lambda}(n_x-1)
 +\hbar\Omega\sum_{x\in\Lambda}b_x^\dagger b_x,
\]
and the hopping is dressed by displacement operators,
\[
 V_{\mathrm{LF}}=-t\sum_{e\in E(\Lambda)}\widetilde B_e,\qquad
 \widetilde B_{\{x,y\}}=\sum_{\sigma}\bigl(D_x^\dagger D_y\,c_{x\sigma}^\dagger
 c_{y\sigma}+D_y^\dagger D_x\,c_{y\sigma}^\dagger c_{x\sigma}\bigr),
\qquad
 D_x=\exp\!\bigl(-\alpha_{\mathrm{LF}}(b_x^\dagger-b_x)\bigr).
\]
Since $\Dop_x$ carries charge two, $\mathcal U\Dop_x\mathcal U^\dagger
=D_x^2\Dop_x$, and $\|D_x^2\|_{\mathrm{op}}=1$.

Expand $e^{-\beta H_{\mathrm{LF}}}$ in a time-ordered Dyson series in the
interaction picture of $H_0$ and insert, along the imaginary-time circle
$[0,\beta)$, the resolution of the occupation basis
\begin{equation}\label{eq:path}
 |n(\tau)\rangle=\bigotimes_{x\in\Lambda}
 |n_{x\uparrow}(\tau),\,n_{x\downarrow}(\tau)\rangle .
\end{equation}
The oscillator trace factorizes over sites. At a single site $x$, the polaronic
counterterm of $H_0$ and the two-time Gaussian contraction,
\begin{equation}\label{eq:contraction}
 \frac U2\int_0^\beta\!\bigl(n_x(\tau)-1\bigr)^2d\tau
 -\frac12\int_0^\beta\!\!\int_0^\beta d\tau\,d\tau'\,
 \bigl(n_x(\tau)-1\bigr)K_{\beta,\Omega}(\tau-\tau')
 \bigl(n_x(\tau')-1\bigr),
\end{equation}
together with $\int_0^\beta K_{\beta,\Omega}(u)\,du=U$, leave the residual weight
\begin{equation}\label{eq:phi}
 \Phi_{\beta,\Omega}[n_x]=\frac14\int_0^\beta\!\!\!\int_0^\beta
 \bigl(n_x(\tau)-n_x(\tau')\bigr)^2K_{\beta,\Omega}(\tau-\tau')\,d\tau\,d\tau'
 =\frac{U}{2\beta}\sum_{m\ne0}\frac{(2\pi m)^2}{(\beta\Omega)^2+(2\pi m)^2}
 \bigl|\hat n_{x,m}\bigr|^2\ge0,
\end{equation}
with the positive periodic kernel
\begin{equation}\label{eq:kernel}
\begin{gathered}
 K_{\beta,\Omega}(u)=\frac{U\Omega}{2}\,
 \frac{\cosh\bigl(\Omega(\beta/2-|u|)\bigr)}{\sinh(\beta\Omega/2)}
 =\frac U\beta\sum_{m\in\mathbb{Z}}
 \frac{(\beta\Omega)^2}{(\beta\Omega)^2+(2\pi m)^2}\,e^{2\pi imu/\beta}>0,\\
 \int_0^\beta K_{\beta,\Omega}(u)\,du=U .
\end{gathered}
\end{equation}
The normalized free-oscillator expectation of the time-ordered displacement string
at site $x$ equals $\exp(-\Phi_{\beta,\Omega}[n_x])$, with or without the equal-time
insertions $D_x^2\Dop_x$ and $(D_y^\dagger)^2\Dop_y^\dagger$ at $\tau=0$, which are
absorbed as jumps of the path. On a constant path $n_x(\tau)\equiv n_x(0)$ the
difference in \eqref{eq:phi} vanishes and $\exp(-\Phi_{\beta,\Omega}[n_x])=1$; the
zero Fourier mode contributes nothing. Because $K_{\beta,\Omega}>0$ and the squared
difference is pointwise nonnegative, every path obeys
\begin{equation}\label{eq:contractionbound}
 \exp\!\left(-\sum_{x\in\Lambda}\Phi_{\beta,\Omega}[n_x]\right)\le1
 \qquad\text{for every }\beta>0,\ \Omega>0,\ M\ge4 .
\end{equation}
The raw Gaussian weight obtained by integrating out the oscillators is not bounded
by one; it is the residual \eqref{eq:phi} after the atomic counterterm that is
bounded, and that is the quantity used below. The step is local in the site index
and independent of $d$.

Two remarks are needed for what follows. First, the bound \eqref{eq:contractionbound} is
uniform as $\Omega\downarrow0$ through positive values, although $\alpha_{\mathrm{LF}}$
diverges there, because $\alpha_{\mathrm{LF}}$ does not enter \eqref{eq:phi}; at
fixed $\beta$ the kernel tends to $U/\beta$. No assertion is made at $\Omega=0$.
Second, a closed path has neutral jumps at each site, and the displacement trace
factorizes over sites, so every oscillator cross-action remains inside the connected
spatial support on which it occurs. No factorization in time is used anywhere.

\section{Signless majorant and connected supports}\label{sec:supports}

In the occupation basis each undirected bond hop $B_e$, $e=\{u,v\}$, has matrix
elements $\pm1$ whose fermionic signs have modulus one. Replacing every matrix
element by its absolute value replaces $B_e$ by the signless two-site bond operator,
again written $B_e$, acting on the sixteen-dimensional bond Fock space $\mathcal F_e$;
a local gauge choice removes signs on a single bond.

Introduce the bond-usage scale $L=L(d)>1$ and write each bond weight as
$t=(t/L)L$. Consider the Dyson expansion of a source-perturbed partition function
$Z_\Lambda(\xi;r)$, defined in Section~\ref{sec:kp}, in which the two-site
pair-transfer labels $E_{\mathrm{src}}=\{\{x,x+r\}:x\in\Lambda\}$ carry weight
$\beta L|\xi|$ each, and $E_{\mathrm{src}}=\emptyset$ when $\xi=0$. A term of order
$m$ specifies labelled bond events $(e_1,\dots,e_m)$ at ordered times
$0<\tau_1<\dots<\tau_m<\beta$.

\paragraph{Two groupings.}
Let $\widetilde E=\{e_1,\dots,e_m\}$ be the set of distinct active labels.

(i) \emph{Connected bond sets.} Decompose $\widetilde E$ into bond components, two
bonds being adjacent when they share an endpoint, and give a component $E_j$ with
$k_j=|E_j|$ distinct bonds the factor $L^{-k_j}$, using $L^{-1}\le1$ on the first
appearance of each distinct bond and $1$ afterwards. Summing over all time-ordered
sequences whose distinct labels are $E_j$, with the covering restriction dropped,
reconstructs the exponential $\exp(\beta t\sum_{e\in E_j}B_e)$ on the vertex span
$V(E_j)$. This is the route of Propositions~\ref{prop:a0} and~\ref{prop:b1}.

(ii) \emph{Connected vertex supports.} Let $S_1\sqcup\dots\sqcup S_p$ be the vertex
sets of the connected components of
$(\bigcup_{e\in\widetilde E}e,\widetilde E)$ in the source-extended graph
$\Lambda_d^+=(\Lambda,E(\Lambda)\cup E_{\mathrm{src}})$. Each $S_j$ is connected and
has $v_j=|S_j|\ge2$, so any spanning tree of $(S_j,\widetilde E|_{S_j})$ has
$v_j-1$ edges and
\begin{equation}\label{eq:handshakecount}
 k_j=\bigl|\widetilde E|_{S_j}\bigr|\ge v_j-1\ge\frac{v_j}{2},
 \qquad\text{hence}\qquad L^{-k_j}\le L^{-(v_j-1)} .
\end{equation}
The map $(e_1,\dots,e_m)\mapsto(S_1,\dots,S_p)$ partitions the active sequences into
classes indexed by pairwise disjoint connected vertex sets. Within one class, summing
over all sequences whose active labels lie in the induced label set
$E_{\Lambda_d^+}(S_j)=\{e\in E(\Lambda)\cup E_{\mathrm{src}}:e\subset S_j\}$, with
both the spanning and the covering restrictions dropped, is an upper bound for the
sum of positive trajectory weights supported on $S_j$; it equals a single
exponential trace on the Fock space of $S_j$ with all induced labels present. Each
connected vertex support therefore contributes one polymer activity, instead of the
$2^{|E(S)|}$ edge subsets or the Euler tours of each subset that the bond-set route
must resum.

The two ingredients that decide the final constant are then the size of that
activity, which is governed by a quadratic-form estimate on the induced subgraph of
$S$, and the number of connected vertex supports containing a given site, which is
governed by animal counting. Sections~\ref{sec:forms} and~\ref{sec:animals} treat
them in turn.

\section{Quadratic-form estimates on induced subgraphs}\label{sec:forms}

Let $S\subset\Lambda$ be a finite vertex set with $|S|=v\ge2$ and let
$E_{\mathrm{phys}}(S)=\{e\in E(\Lambda):e\subset S\}$. Let
$n_{\mathrm{odd},x}=\mathbf 1_{\{n_x=1\}}\in\{0,1\}$ be the projection onto odd site
occupation, so that $(n_x-1)^2=1-n_{\mathrm{odd},x}$, and set
\[
 \Nodd(S)=\sum_{x\in S}n_{\mathrm{odd},x},\qquad
 N_{\mathrm{odd},e}=n_{\mathrm{odd},u}+n_{\mathrm{odd},v}\quad(e=\{u,v\}).
\]
Every $x\in S$ is incident to at most $\Dphys(d)=2d$ bonds of
$E_{\mathrm{phys}}(S)$, so summing $N_{\mathrm{odd},e}$ over edges gives the exact
graph-degree handshake identity
\begin{equation}\label{eq:handshake}
 \sum_{e\in E_{\mathrm{phys}}(S)}N_{\mathrm{odd},e}
 =\sum_{x\in S}\deg_{E_{\mathrm{phys}}(S)}(x)\,n_{\mathrm{odd},x}
 \le 2d\,\Nodd(S),
\end{equation}
operator-wise, since the projections $n_{\mathrm{odd},x}$ are positive
semidefinite and commute.

\subsection{Edge-local sum of squares, unconditional in the hopping}

The odd-site gap of the atomic Hamiltonian at site $x$ is $\frac U2$ on the
singleton sector, and it pays for the hopping of a single bond through the
following estimate on the sixteen-dimensional bond Fock space $\mathcal F_e$.

\begin{lemma}[Edge-local sum of squares]\label{lem:sos}
For every bond $e$, every $g\ge0$ and every $\alpha>0$,
\begin{equation}\label{eq:localform}
 g\,B_e\ \preceq\ \alpha\,N_{\mathrm{odd},e}+\frac{2g^2}{\alpha}\,I_{16},
\end{equation}
equivalently $R_e(g,\alpha):=\alpha N_{\mathrm{odd},e}+\frac{2g^2}{\alpha}I_{16}-gB_e
\succeq0$, with no upper bound on $g$.
\end{lemma}

\begin{proof}
$B_e$ and $N_{\mathrm{odd},e}$ conserve the bond particle number
$N_e=n_u+n_v\in\{0,1,2,3,4\}$ and are graded by it, so the sectors are invariant.
They are also invariant under the spin rotation that swaps the two sites together
with the spin, so it suffices to treat each sector once.

\emph{$N_e\in\{0,4\}$.} Here $B_e=0$ and $N_{\mathrm{odd},e}=0$, so
$R_e=\frac{2g^2}{\alpha}I\succeq0$.

\emph{$N_e\in\{1,3\}$.} Here $N_{\mathrm{odd},e}=I_4$ and $B_e$ splits into two
decoupled spin doublets with matrix $\begin{pmatrix}0&1\\1&0\end{pmatrix}$. For each
doublet and $\psi=(\psi_1,\psi_2)$,
\[
 \langle\psi,R_e\psi\rangle
 =\Bigl|\sqrt\alpha\,\psi_1-\frac{g}{2\sqrt\alpha}\psi_2\Bigr|^2
 +\Bigl|\sqrt\alpha\,\psi_2-\frac{g}{2\sqrt\alpha}\psi_1\Bigr|^2
 +\frac{7g^2}{4\alpha}\bigl(|\psi_1|^2+|\psi_2|^2\bigr)\ge0 ,
\]
This is the sector in which the odd-preserving hop lives, and it is here that the
completion of the square absorbs it for every $g$, not merely for small $g/U$.

\emph{$N_e=2$ ($\dim=6$).} The four odd-odd states have $N_{\mathrm{odd},e}=2$ and
the two doubly-occupied states have $N_{\mathrm{odd},e}=0$. On the symmetric
combinations
\[
 |o_+\rangle=\tfrac1{\sqrt2}\bigl(|1,0,0,1\rangle+|0,1,1,0\rangle\bigr),\qquad
 |e_+\rangle=\tfrac1{\sqrt2}\bigl(|1,1,0,0\rangle+|0,0,1,1\rangle\bigr),
\]
the signless operator acts as $B_e|e_+\rangle=2|o_+\rangle$ and
$B_e|o_+\rangle=2|e_+\rangle$, and it annihilates the orthogonal four-dimensional
subspace. Hence
\[
 R_e\big|_{\{|e_+\rangle,|o_+\rangle\}}
 =\begin{pmatrix}\frac{2g^2}{\alpha}&-2g\\-2g&2\alpha+\frac{2g^2}{\alpha}\end{pmatrix}
 =\begin{pmatrix}-\frac{\sqrt2 g}{\sqrt\alpha}\\ \sqrt{2\alpha}\end{pmatrix}
  \begin{pmatrix}-\frac{\sqrt2 g}{\sqrt\alpha}&\sqrt{2\alpha}\end{pmatrix}
 +\begin{pmatrix}0&0\\0&\frac{2g^2}{\alpha}\end{pmatrix}\succeq0 .
\]
This completes Lemma~\ref{lem:sos}.
\end{proof}

Summing \eqref{eq:localform} over $e\in E_{\mathrm{phys}}(S)$ with
$g=L(d)\,t$ and applying \eqref{eq:handshake} gives, for every $g\ge0$ and every
$\alpha>0$,
\begin{equation}\label{eq:sumform}
 g\!\!\sum_{e\in E_{\mathrm{phys}}(S)}\!\!B_e
 \le\alpha\!\!\sum_{e\in E_{\mathrm{phys}}(S)}\!\!N_{\mathrm{odd},e}
  +\frac{2g^2}{\alpha}\bigl|E_{\mathrm{phys}}(S)\bigr|I
 \le2d\alpha\,\Nodd(S)+\frac{2g^2}{\alpha}\bigl|E_{\mathrm{phys}}(S)\bigr|I .
\end{equation}
The gap share is $\frac U4\Nodd(S)$ in both the unconditional and the conditional
route, leaving the remaining $\frac U4\Nodd(S)$ of the atomic odd-site penalty
$\frac U2\Nodd(S)$ for the suppression of odd-site occupancy. Setting
\begin{equation}\label{eq:alpha}
 2d\alpha=\frac U4\ \Longleftrightarrow\ \alpha=\frac{U}{8d}
 \quad\Longrightarrow\quad \frac{2g^2}{\alpha}=\frac{16\,d\,g^2}{U},
\end{equation}
and using $|E_{\mathrm{phys}}(S)|\le\frac12\Dphys(d)\,|S|=d\,|S|$, \eqref{eq:sumform}
becomes the unconditional form bound
\begin{equation}\label{eq:unconditional}
 g\!\!\sum_{e\in E_{\mathrm{phys}}(S)}\!\!B_e
 \le\frac U4\,\Nodd(S)+\frac{16\,d\,g^2}{U}\bigl|E_{\mathrm{phys}}(S)\bigr|I
 \le\frac U4\,\Nodd(S)+\frac{16\,d^2g^2}{U}\,|S|\,I
 \qquad(\forall g\ge0).
\end{equation}
Because \eqref{eq:localform} holds for all $g\ge0$, Theorem~\ref{thm:main} and
Proposition~\ref{prop:a0} carry no upper bound on $t/U$.

\subsection{The conditional global form and its smallness cutoff}

The alternative route decomposes $B_{\mathrm{phys}}(S)=\sum_{e}B_e$ globally on
$\mathcal F_S$ according to the change of $\Nodd$. Each of the $k=|E_{\mathrm{phys}}(S)|$
bonds supports at most two spin-creation transitions and each odd site at most
$\Dphys(d)=2d$ annihilation transitions, so weighted Cauchy-Schwarz with parameter
$h>0$ gives
\begin{equation}\label{eq:cs}
\begin{gathered}
 g(B_++B_-)\le2kgh\,I+\frac{2d\,g}{h}\,\Nodd(S),
 \qquad h=\frac{16dg}{U},\\
\Longrightarrow\quad
 g(B_++B_-)\le\frac{32\,d\,k\,g^2}{U}\,I+\frac U8\,\Nodd(S).
\end{gathered}
\end{equation}
Each odd-preserving hop touches an odd site, so each configuration has at most
$2\Dphys(d)\Nodd(S)=4d\,\Nodd(S)$ such transitions, and
\begin{equation}\label{eq:b0}
 gB_0\le2\Dphys(d)\,g\,\Nodd(S)=4dg\,\Nodd(S).
\end{equation}
Absorbing \eqref{eq:b0} into the remaining $\frac U8\Nodd(S)$ requires
\begin{equation}\label{eq:smallcutoff}
 2\Dphys(d)\,g\le\frac U8
 \quad\Longleftrightarrow\quad
 \frac{g}{U}\le\frac{1}{16\,\Dphys(d)}=\frac{1}{32\,d},
\end{equation}
which with $g=Lt$ is the cutoff $t/U\le1/(32\,d\,L)$; applied with the two scales
of Propositions~\ref{prop:b1} and~\ref{prop:b2} it gives the thresholds
$\varepsilon_0$ stated there. The condition is not an artefact of the
estimates: in the $\Delta\Nodd=0$ channel the global count discards the local
information that these hops act only within single bonds. On a single bond,
$\Delta\Nodd=0$ transitions occur exclusively in the sectors $N_e\in\{1,3\}$ where
$N_{\mathrm{odd},e}=1$, and there \eqref{eq:localform} already absorbs $gB_{0,e}$ for
all $g\ge0$. Removing the cutoff is therefore not a matter of a better constant but of
using the bond-local structure of the kinetic energy, and that is what
Lemma~\ref{lem:sos} does.

\subsection{Monotonicity and the single-site atomic fugacity ratio}

\begin{lemma}[Trace monotonicity]\label{lem:monotone}
Let $X\preceq Y$ be self-adjoint operators on a finite-dimensional space. Then all
eigenvalues ordered increasingly satisfy $\lambda_j(X)\le\lambda_j(Y)$, and
$\Tr e^{\beta X}\le\Tr e^{\beta Y}$ for every $\beta>0$.
\end{lemma}

\begin{proof}
For any fixed $j$, the min-max formula
$\lambda_j(Y)=\max_{\dim\mathcal V=j}\min_{v\in\mathcal V,\,\|v\|=1}\langle v,Yv\rangle$
and the pointwise inequality $\langle v,Xv\rangle\le\langle v,Yv\rangle$ give
$\lambda_j(X)\le\lambda_j(Y)$ for every $j$. As $s\mapsto e^{\beta s}$ is
strictly increasing,
$\Tr e^{\beta X}=\sum_je^{\beta\lambda_j(X)}\le\sum_je^{\beta\lambda_j(Y)}=\Tr e^{\beta Y}$.
\end{proof}

On $\mathcal F_S$ define
$H_{\mathrm{at}}(S)=-\frac U2\sum_{x\in S}(n_x-1)^2-\mu\sum_{x\in S}(n_x-1)
 =-\frac U2v\,I+\frac U2\Nodd(S)-\mu\sum_{x\in S}(n_x-1)$.
Taking $X=-H_{\mathrm{at}}(S)+gB_{\mathrm{phys}}(S)$ and
$Y=-H_{\mathrm{at}}(S)+\frac U4\Nodd(S)+\frac{16d^2L(d)^2t^2}{U}v\,I$, which dominates
$X$ by \eqref{eq:unconditional}, and unwrapping the atomic identity
$\frac U2\Nodd-\frac U4\Nodd$ gives
\begin{equation}\label{eq:tracedrop}
 \frac{\Tr_{\mathcal F_S}\bigl(e^{-\beta H_{\mathrm{at}}(S)+\beta L(d)t\,B_{\mathrm{phys}}(S)}\bigr)}
      {\Tr_{\mathcal F_S}\bigl(e^{-\beta H_{\mathrm{at}}(S)}\bigr)}
 \le e^{\delta_{\mathrm{phys}}(d)\,v}\prod_{x\in S}\frac{z_b(\mu)}{z_a(\mu)},
\end{equation}
where
\begin{equation}\label{eq:delta}
 \delta_{\mathrm{phys}}(d)=\frac{16\,d^2L(d)^2\beta t^2}{U}
 \le\frac{16\,d^2L(d)^2}{\frac{256}{3}d^2L(d)^2}=\frac{3}{16}<\frac14 ,
\end{equation}
and, after factoring out $e^{\beta U/2-\beta\mu}$ and with $z=e^{\beta\mu}>0$ and
$w=e^{-\beta U/4}\in(0,1)$,
\begin{equation}\label{eq:zazb}
 z_a(\mu)=1+2ze^{-\beta U/2}+z^2=1+2zw^2+z^2,\qquad
 z_b(\mu)=1+2ze^{-\beta U/4}+z^2=1+2zw+z^2 .
\end{equation}
The bivariate polynomial identity
\begin{equation}\label{eq:fivezb}
 5z_a(\mu)-4z_b(\mu)=(1-z)^2+2z(2w-1)^2+2zw^2>0
 \qquad\Longrightarrow\qquad\frac{z_b(\mu)}{z_a(\mu)}<\frac54 ,
\end{equation}
which holds for every $z>0$ and every $w\in(0,1)$, together with
\eqref{eq:tracedrop} and \eqref{eq:delta}, bounds the normalized physical trace on
$S$ by $\bigl(\frac54e^{\delta_{\mathrm{phys}}(d)}\bigr)^v$ with
$\delta_{\mathrm{phys}}(d)\le\frac{3}{16}<\frac14$, uniformly in $\mu\in\mathbb R$,
$\Omega>0$ and $M\ge4$. Nothing in \eqref{eq:tracedrop} depends on the dimension
beyond the explicit occurrences of $d$.

\section{Counting connected supports by Bethe-tree animals}\label{sec:animals}

\begin{lemma}[Bethe-tree injection and exact Lagrange formula]\label{lem:animals}
Let $G$ be a simple graph of maximum degree $\Delta\ge2$, let
$\mathcal A_G(x_0,v)$ be the set of connected $v$-element vertex sets containing
$x_0$, and let $a_v(G)=\max_{x_0}|\mathcal A_G(x_0,v)|$. Let $\mathbb T_\Delta$ be
the infinite $\Delta$-regular tree rooted at $o$, whose root has $\Delta$ children and
every other vertex $\Delta-1$, and let $T_\Delta(v)$ count its connected $v$-element
vertex sets containing $o$. Then:
\begin{enumerate}
\item there is an explicit injection $\mathcal A_G(x_0,v)\hookrightarrow
 \mathcal A_{\mathbb T_\Delta}(o,v)$, so $a_v(G)\le T_\Delta(v)$;
\item the generating function $R_\Delta(x)=\sum_{v\ge1}T_\Delta(v)x^v$ satisfies the
 exact functional equation
 \begin{equation}\label{eq:functional}
 R_\Delta(x)=x\bigl(1+f_\Delta(x)\bigr)^\Delta,\qquad
 f_\Delta(x)=x\bigl(1+f_\Delta(x)\bigr)^{\Delta-1},
 \end{equation}
 and by Lagrange inversion
\begin{equation}\label{eq:lagrange}
 T_\Delta(1)=1,\qquad T_\Delta(v)=\frac{\Delta}{v-1}\binom{v(\Delta-1)}{v-2}
 \quad(v\ge2);
\end{equation}
\item on $\mathbb Z^d$ with $\Delta=2d$, $a_v(\mathbb Z^d)=T_{2d}(v)$ for
 $v\in\{1,2,3\}$, and at $v=4$ the exact plaquette identity holds:
\begin{equation}\label{eq:plaquette}
 T_{2d}(4)-a_4(\mathbb Z^d)=12\binom d2=6d(d-1);
\end{equation}
\item for every $\Delta\ge2$ and $v\ge1$, $T_\Delta(v)\le(e\Delta)^{v-1}$.
\end{enumerate}
\end{lemma}

\begin{proof}
(1) Fix a total order on $V(G)$ and injectively label, at each vertex, the incident
edges by $1,\dots,\Delta$. Breadth-first search in $G[S]$ from $x_0$, visiting
neighbours in increasing order, produces a unique rooted spanning tree $\tau(S)$ of
$S$. Lifting $\tau(S)$ to $\mathbb T_\Delta$ along the edge labels, with
$x_0\mapsto o$, gives a rooted subtree of size $v$: the walk is self-avoiding
because $\tau(S)$ is a tree, and each non-root vertex uses at most $\Delta-1$
directions other than its parent. Distinct $S$ give distinct trees, since $S$ is the
vertex set of $\iota(S)$.

(2) At the root, each of the $\Delta$ branches independently carries either nothing
or a rooted subtree of the $(\Delta-1)$-ary branch, whence
$R_\Delta=x(1+f_\Delta)^\Delta$ and $f_\Delta=x(1+f_\Delta)^{\Delta-1}$. Writing
$R_\Delta=x\,g(f_\Delta)$ with $g(u)=(1+u)^\Delta$ and
$\phi(u)=(1+u)^{\Delta-1}$, Lagrange inversion with $n=v-1\ge1$ gives
\[
 T_\Delta(v)=[x^n]g(f_\Delta(x))=\frac1n[u^{n-1}]\,g'(u)\phi(u)^n
 =\frac1{v-1}[u^{v-2}]\,\Delta(1+u)^{v(\Delta-1)}
 =\frac{\Delta}{v-1}\binom{v(\Delta-1)}{v-2}.
\]

(3) $\mathbb Z^d$ is bipartite, so no induced subgraph on at most three vertices has a
cycle and no two distinct rooted trees on at most three vertices can cover the same
vertex set; hence equality for $v\le3$. On $\mathbb Z^d$ the only cycles of length at
most four containing the origin are the $4\binom d2$ coordinate plaquettes of side
two. Each plaquette has four spanning trees, which lift to four distinct rooted
subtrees of $\mathbb T_{2d}$ while giving a single four-vertex animal of
$\mathbb Z^d$, so the excess is exactly $(4-1)\cdot4\binom d2=12\binom d2$.

(4) $T_\Delta(1)=1=(e\Delta)^0$ and $T_\Delta(2)=\Delta<e\Delta$. For $v\ge3$, use
$\binom Nm\le N^m/m!$ with $N=v(\Delta-1)$ and $m=v-2\ge1$, and the bound
$v^{v-2}/(v-1)!=v^{v-1}/v!\le e^v/v<e^{v-1}$, which follows from
$e^v=\sum_jv^j/j!>v^v/v!$ and $v\ge3>e$:
\[
 T_\Delta(v)=\frac{\Delta}{v-1}\binom{v(\Delta-1)}{v-2}
 \le\Delta(\Delta-1)^{v-2}\frac{v^{v-1}}{v!}<\Delta^{v-1}e^{v-1}=(e\Delta)^{v-1}.
\]
\end{proof}

The encoding of animals by rooted trees goes back to Klarner \cite{klarner}; the
proof above is self-contained and does not use any external growth constant.
Part (3) is the reason a four-vertex animal can hide a cycle, and it is the only place
where the dimension enters the animal count at low order: the plaquette excess
$12\binom{d}{2}$ grows like $6d^2$, while
$T_{2d}(4)=\frac{2d}{3}\binom{8d-4}{2}$ grows like $\frac{64}{3}d^3$.

Rather than use the coarse bound (4) inside the Koteck\'y-Preiss sum, we evaluate the
exact generating function at one rational point.

\begin{lemma}[Generating function at $y_0$]\label{lem:gf}
For every integer $\Delta\ge2$ put
\begin{equation}\label{eq:y0}
 y_0=\frac{1}{6\Delta},\qquad x_0=\frac{y_0}{(1+y_0)^{\Delta-1}} .
\end{equation}
Then for every $x\in[0,x_0]$ the series $R_\Delta(x)$ converges and
\begin{equation}\label{eq:gfbound}
 \sum_{v\ge2}T_\Delta(v)\,x^v\le x\bigl((1+y_0)^\Delta-1\bigr)
 <x\bigl(e^{1/6}-1\bigr)<\frac x5 .
\end{equation}
\end{lemma}

\begin{proof}
$f_\Delta$ has nonnegative coefficients and $f_\Delta(x_0)=y_0$, so for
$x\in[0,x_0]$ we have $f_\Delta(x)/x=(1+f_\Delta(x))^{\Delta-1}\le(1+y_0)^{\Delta-1}$,
and monotone convergence of the coefficients gives
\[
 \sum_{v\ge2}T_\Delta(v)x^v
 \le x\sum_{v\ge2}T_\Delta(v)x_0^{v-1}
 =x\Bigl(\frac{R_\Delta(x_0)}{x_0}-1\Bigr)=x\bigl((1+y_0)^\Delta-1\bigr).
\]
By the binomial series $(1+\frac1{6\Delta})^\Delta<e^{1/6}$, and by the rational
Taylor enclosure
\begin{equation}\label{eq:eenc}
 e=\sum_{n\ge0}\frac1{n!}
 <\sum_{n=0}^5\frac1{n!}+\frac1{6!}\sum_{m\ge0}\Bigl(\frac17\Bigr)^m
 =\frac{163}{60}+\frac1{720}\cdot\frac76=\frac{11743}{4320}<\frac{1175}{432}
 <\Bigl(\frac65\Bigr)^6=\frac{46656}{15625},
\end{equation}
so that $e^{1/6}<6/5$ and $(1+y_0)^\Delta-1<1/5$.
\end{proof}

\section{Polymer activities, sources, and Koteck\'y-Preiss convergence}\label{sec:kp}

\subsection{The pair-transfer source and the augmented degree}

Fix a nonzero separation $r\in\Lambda\setminus\{0\}$ with $|r|_1\ge1$ and a real
phase $\theta$. The number-conserving pair-transfer source on the pair
$e=\{x,x+r\}$ is
\begin{equation}\label{eq:source}
 X_e(\theta)=e^{i\theta}\Dop_x^\dagger\Dop_{x+r}+e^{-i\theta}\Dop_{x+r}^\dagger\Dop_x,
 \qquad\|X_e(\theta)\|_{\mathrm{op}}\le2,
\end{equation}
and the source-perturbed partition function is
\[
 Z_\Lambda(\xi;r)=\Tr\Bigl[e^{-\beta H_\Lambda(\mu)}\,
 \mathcal T_\tau\exp\Bigl(\int_0^\beta\!d\tau\;\xi e^{i\theta}
 \sum_{x\in\Lambda}\Dop_x^\dagger\Dop_{x+r}(\tau)+\text{h.c.}\Bigr)\Bigr],
 \qquad \xi\in\mathbb C,
\]
where the fields are Heisenberg evolved and $\mathcal T_\tau$ is imaginary-time
ordering. Writing $p_\Lambda(\xi)=V^{-1}\log Z_\Lambda(\xi;r)$,
\[
 p_\Lambda'(0)=\frac{2\beta}{V}\sum_{x\in\Lambda}
 \operatorname{Re}\bigl(e^{i\theta}\langle\Dop_x^\dagger\Dop_{x+r}\rangle\bigr),
\]
so a bound at fixed $r$ gives the correlation function at every separation, and the
translation average is what enters the canonical transfer of
Section~\ref{sec:canonical}.

For each source pair, the operator $X_e(\theta)$ transfers a pair between the two
sites, mapping $|2,0\rangle\leftrightarrow|0,2\rangle$, a state with one site doubly
occupied and the other empty; on both active local configurations $n_u,n_v\in\{0,2\}$,
so $N_{\mathrm{odd},e}=0$. The odd-occupation
penalty therefore vanishes on the active subspace of $X_e$ and cannot absorb it
through a completion of squares. What replaces it is a norm bound. A source pair
contributes at most one term per site of $S$, and each site of $\Lambda$ is the
bottom end of at most one source pair, so a connected vertex support $S$ contains at
most $v$ source pairs. Hence, on the source disk
\begin{equation}\label{eq:disk}
 |\xi|\le\frac{\tau}{16L(d)}=\frac{1}{16\beta L(d)},
\end{equation}
the source perturbation of a support $S$ obeys, by Lemma~\ref{lem:monotone} and the
size estimate,
\begin{equation}\label{eq:sourcepenalty}
 \beta L(d)|\xi|\cdot v\cdot\|X_e\|_{\mathrm{op}}\le2\beta L(d)|\xi|\,v\le\frac v8<\frac v4 .
\end{equation}
Adding at most two source edges per site at displacements $\pm r$ raises the maximum
degree of the underlying simple graph from $\Dphys(d)=2d$ to
\begin{equation}\label{eq:daug}
 \Daug(d)=\Dphys(d)+2=2d+2 .
\end{equation}
If a source label shares its endpoints with a nearest-neighbour bond or with the
reverse translate, the resulting multigraph has parallel edges; connectedness of $S$
depends only on the underlying simple graph of maximum degree $\Daug(d)$, so
Lemma~\ref{lem:animals} applies unchanged. In the translation-averaged source each
site initiates at most one term with both ends in $S$, so again at most $v$ source
pairs per support.

\subsection{Activity bound}

For a connected vertex support $S$ with $v=|S|\ge2$ in the source-extended graph,
combine the first-appearance scaling factor $L(d)^{-(v-1)}$ from
\eqref{eq:handshakecount}, the single-site atomic ratio $(z_b/z_a)^v<(5/4)^v$ from
\eqref{eq:fivezb}, the physical trace penalty $e^{\delta_{\mathrm{phys}}(d)v}$ from
\eqref{eq:delta}, and the scalar source penalty $e^{2\beta L(d)|\xi|v}\le e^{v/8}$ from
\eqref{eq:sourcepenalty}. The activity is then bounded by
\begin{equation}\label{eq:activity}
 |w(S)|\le L(d)^{-(v-1)}\Bigl(\frac54\,e^{\delta_{\mathrm{tot}}(d)}\Bigr)^{v},
 \qquad
 \delta_{\mathrm{tot}}(d)=\delta_{\mathrm{phys}}(d)+\begin{cases}0,& \xi=0,\\
 2\beta L(d)|\xi|,& \xi\ne0.\end{cases}
\end{equation}
In the physical ensemble $\delta_{\mathrm{tot}}=\delta_{\mathrm{phys}}\le\frac{3}{16}<\frac14$,
and since $(4/3)^4=\frac{256}{81}>\frac{1175}{432}>e$ we have
$e^{\delta_{\mathrm{phys}}}<e^{1/4}<\frac43$, so that
\begin{equation}\label{eq:xphys}
 x_{\mathrm{phys}}:=\frac{5}{4L(d)}\,e^{\delta_{\mathrm{phys}}(d)}e,
 \qquad
 L(d)\,x_{\mathrm{phys}}<\frac54\cdot\frac43\cdot\frac{1175}{432}
 =\frac{5875}{1296}\approx4.5332<5 .
\end{equation}
In the source-perturbed ensemble on the disk \eqref{eq:disk},
\begin{equation}\label{eq:deltatot}
 \delta_{\mathrm{tot}}(d)\le\frac{3}{16}+\frac18=\frac{5}{16},
\end{equation}
and the exact rational enclosure
\[
 \Bigl(\frac{11}{8}\Bigr)^{16}=\frac{45\,949\,729\,863\,572\,161}{281\,474\,976\,710\,656}
 >\Bigl(\frac{1175}{432}\Bigr)^{5}=\frac{2\,239\,697\,333\,984\,375}{15\,045\,919\,506\,432}
 >e^5
 \ \Longrightarrow\ e^{5/16}<\frac{11}{8}
\]
gives
\begin{equation}\label{eq:xsrc}
 x_{\mathrm{src}}:=\frac{5}{4L(d)}\,e^{\delta_{\mathrm{tot}}(d)}e,
 \qquad
 L(d)\,x_{\mathrm{src}}\le\frac54e^{5/16}e<\frac54\cdot\frac{11}{8}\cdot\frac{1175}{432}
 =\frac{64\,625}{13\,824}\approx4.6748<5 .
\end{equation}

\subsection{Koteck\'y-Preiss criterion with sources}

\begin{lemma}[Polymer convergence and support deletion]\label{lem:kp}
Let $\mathcal S_\Lambda$ be a family of connected vertex supports $S\subset\Lambda$ with
$|S|\ge2$, let $w:\mathcal S_\Lambda\to\mathbb C$ be activities, and declare
$S,S'$ compatible iff $S\cap S'=\emptyset$. Set $a(S)=|S|$ and $d(S)=0$ and suppose
\begin{equation}\label{eq:kpcond}
 \sup_{x_0\in\Lambda}\sum_{S\in\mathcal S_\Lambda,\,S\ni x_0}|w(S)|\,e^{|S|}\le1 .
\end{equation}
Then:
\begin{enumerate}
\item the cluster expansion of $\log\sum_{\Gamma\text{ compatible}}\prod_{S\in\Gamma}w(S)$
converges absolutely and its value is analytic in the activities on the polydisc
defined by \eqref{eq:kpcond}, uniformly in $\Lambda$;
\item for every multi-index derivative the corresponding truncated sum is bounded by
the corresponding derivative of $\prod(1+|w(S)|e^{|S|})$ at the same point, so in
particular $|\partial_\xi\log Z_\Lambda(\xi;r)|$ at $\xi=0$ is bounded by the
linear activity sum;
\item for every $S_0\in\mathcal S_\Lambda$ the ratio of polymer partition functions
obtained by deleting $S_0$ obeys
\begin{equation}\label{eq:deletion}
 \Bigl|\frac{Z_{\Lambda\setminus S_0}}{Z_\Lambda}\Bigr|\le e^{|S_0|},
\end{equation}
and the same bound holds for the source-perturbed ensembles on the disk
\eqref{eq:disk}.
\end{enumerate}
\end{lemma}

\begin{proof}
Statement (1) is Theorem~1 of Koteck\'y and Preiss \cite{kp} with
$a(S)=|S|$, $d(S)=0$ and compatibility given by vertex disjointness: the condition
$\sum_{S\not\sim S_0}|w(S)|e^{a(S)+d(S)}\le a(S_0)=|S_0|$ reduces, by subadditivity
over the at most $|S_0|$ sites of $S_0$, to \eqref{eq:kpcond}. Statement (2) is the
convergence of the Ursell sums on the polydisc, since convergence at a point of the
polydisc with positive activities gives absolute convergence of every coefficient
series and hence of every derivative. For (3), expand the ratio along the tree-graph
(Penrose) identity: $\log(Z_{\Lambda}/Z_{\Lambda\setminus S_0})$ is a sum over
connected clusters of the gas that meet $S_0$, each cluster contributing a product of
activities whose modulus sum is bounded by $\sum_{x_0\in S_0}1=|S_0|$ because
\eqref{eq:kpcond} is uniform in $\Lambda$. Adding constants to the logarithm gives
\eqref{eq:deletion}. The source-perturbed statements are the same with the activities
of \eqref{eq:activity}, whose pinned sums satisfy the same condition below.
\end{proof}

With $\Daug(d)=2d+2$ and the choice
\begin{equation}\label{eq:L}
 L(d)=36\,\Daug(d)=72(d+1),
\end{equation}
Lemma~\ref{lem:gf} at $y_0=1/(6\Daug(d))$ applies, because
$(1+y_0)^{\Daug-1}<e^{1/6}<6/5$ by \eqref{eq:eenc} and therefore
\begin{equation}\label{eq:x0order}
 x_{\mathrm{phys}}\le x_{\mathrm{src}}<\frac{5}{L(d)}=\frac{5}{36\Daug(d)}
 <x_0=\frac{1}{6\Daug(d)(1+y_0)^{\Daug(d)-1}} .
\end{equation}
Consequently the pinned Koteck\'y-Preiss budgets are strictly below one in every
dimension:
\begin{align}
 \sum_{S\ni x_0,\,|S|\ge2}|w_{\mathrm{phys}}(S)|\,e^{|S|}
 &<\frac{5875}{1296}\Bigl((1+y_0)^{2d+2}-1\Bigr)
 <\frac{5875}{1296}\cdot\frac15=\frac{1175}{1296}<1,\label{eq:budgetphys}\\
 \sum_{S\ni x_0,\,|S|\ge2}|w_{\mathrm{src}}(S)|\,e^{|S|}
 &<\frac{64\,625}{13\,824}\Bigl((1+y_0)^{2d+2}-1\Bigr)
 <\frac{64\,625}{13\,824}\cdot\frac15=\frac{12\,925}{13\,824}<1.\label{eq:budgetsrc}
\end{align}
Equations \eqref{eq:budgetphys} and \eqref{eq:budgetsrc} are the whole reason the
constant $L(d)$ is of order $d$: the animal growth constant of the augmented graph is
$\Daug(d)=2d+2$, so the activity must be paid by a first-appearance factor of order
$d$ per site, and $L$ enters the final constant quadratically.

\section{Decay rate and prefactor}\label{sec:decay}

\subsection{Signed hard-core deletion}

In the physical ensemble $Z_\Lambda$ is a signed hard-core polymer partition
function,
\begin{equation}\label{eq:signed}
 \frac{Z_\Lambda}{Z_{\mathrm{at},\Lambda}}
 =\sum_{\Gamma\text{ compatible}}\ \prod_{S\in\Gamma}w_{\mathrm{signed}}(S),
 \qquad |w_{\mathrm{signed}}(S)|\le|w_{\mathrm{phys}}(S)| ,
\end{equation}
whose activities are bounded in absolute value by the input activities of the
majorant.
The activities are signed because a closed trajectory carries fermionic and
orientational signs; the absolute-value majorant is what \eqref{eq:activity} bounds,
and dropping the alternating terms of the sum is not legitimate. The deletion
bound \eqref{eq:deletion} of Lemma~\ref{lem:kp} is exactly what replaces that
dropping.

Consider distinct $x,y\in\Lambda$ at distance $d_{\mathrm{dist}}=|x-y|_1\ge1$ and the
two-point function $\langle\Dop_x^\dagger\Dop_y\rangle$. Any connected component of
the trajectory support that contains exactly one of the two insertions has zero
trace: it carries a single pair creation or annihilation operator, $\Dop_x^\dagger$ or
$\Dop_y$, which changes the particle number of that component by two, and charge
conservation kills the trace. Hence a single connected support $S_0$ carries both insertions, whence
\begin{equation}\label{eq:size}
 v=|S_0|=k+1\ge d_{\mathrm{dist}}+1,\qquad k:=v-1\ge d_{\mathrm{dist}}\ge1 .
\end{equation}
Since $\|\Dop_x^\dagger\Dop_y\|_{\mathrm{op}}=1$ and the pair insertion preserves the
odd-occupation parity at both endpoints, the marked activity is bounded by
\begin{equation}\label{eq:marked}
 |w_{\mathrm{phys}}(S_0)|\le L(d)^{-k}\Bigl(\frac54e^{\delta_{\mathrm{phys}}(d)}\Bigr)^{k+1}.
\end{equation}
Multiplying \eqref{eq:marked} by the deletion factor $e^{k+1}$ of
\eqref{eq:deletion} and summing over connected supports $S_0\supset\{x,y\}$, whose
number is at most the number $T_{2d}(k+1)$ of size-$(k+1)$ animals rooted at $x$ in a
graph of maximum degree $2d$ (Lemma~\ref{lem:animals}(1)), gives
\begin{equation}\label{eq:geometric}
 \bigl|\langle\Dop_x^\dagger\Dop_y\rangle\bigr|
 \le\sum_{k\ge d_{\mathrm{dist}}}T_{2d}(k+1)L(d)^{-k}
 \Bigl(\frac54e^{\delta_{\mathrm{phys}}(d)}e\Bigr)^{k+1}
 =L(d)\sum_{k\ge d_{\mathrm{dist}}}T_{2d}(k+1)\,x_{\mathrm{phys}}^{k+1}.
\end{equation}

\subsection{Sharp evaluation of the animal generating function}

To make the geometric series uniform rather than coefficient by coefficient, evaluate
$R_{2d}(x)=x(1+f_{2d}(x))^{2d}=\sum_{k\ge0}T_{2d}(k+1)x^{k+1}$ of
Lemma~\ref{lem:animals}(2) at
\begin{equation}\label{eq:y1}
 y_1=\frac1{2d-1},\qquad
 x_1=\frac{y_1}{(1+y_1)^{2d-1}}.
\end{equation}
Since all coefficients are nonnegative and
$\sum_{k\ge0}T_{2d}(k+1)x_1^k=R_{2d}(x_1)/x_1=(1+y_1)^{2d}$, every coefficient obeys
\begin{equation}\label{eq:coeff}
 T_{2d}(k+1)\le\frac{(1+y_1)^{2d}}{x_1^k}\qquad(k\ge0).
\end{equation}
Two estimates remain.

\emph{Step ratio.} For integers $m\ge1$ and $2\le j\le m$,
$\binom mj m^{-j}=\frac1{j!}\prod_{i=1}^{j-1}(1-\frac im)<\frac1{j!}$, so
\begin{equation}\label{eq:elev}
 \Bigl(1+\frac1m\Bigr)^m=2+\sum_{j=2}^m\binom mj\frac1{m^j}
 <\sum_{j=0}^m\frac1{j!}<e<\frac{1175}{432}.
\end{equation}
With $m=2d-1\ge3$ this gives $1/x_1<(2d-1)\cdot\frac{1175}{432}$, and with
$L(d)x_{\mathrm{phys}}<\frac{5875}{1296}$ from \eqref{eq:xphys} and
$L(d)=72(d+1)$,
\begin{equation}\label{eq:qsharp}
 \frac{x_{\mathrm{phys}}}{x_1}
 =\frac{L(d)x_{\mathrm{phys}}}{L(d)x_1}
 <\frac{\frac{5875}{1296}\cdot\frac{1175}{432}(2d-1)}{72(d+1)}
 =\frac{6\,903\,125\,(2d-1)}{40\,310\,784\,(d+1)}
 =q_{\mathrm{sharp}}(d)<q(d)<\frac5{12}<1,
\end{equation}
where the factor $\frac{2d-1}{d+1}>0$ cancels and the middle inequality is the
exact rational comparison
$\frac{6\,903\,125}{40\,310\,784}\approx0.171248<\frac5{24}
=\frac{8\,398\,080}{40\,310\,784}$.

\emph{Prefactor.} The sequence $\bigl(1+\frac1{2d-1}\bigr)^{2d}$ is
$\bigl(1+\frac1n\bigr)^{n+1}$ at $n=2d-1$, which is strictly decreasing in $n$;
it is therefore largest at $d=2$, where
$(1+\frac13)^4=(4/3)^4=\frac{256}{81}$, so for every $d\ge2$
\begin{equation}\label{eq:prefactor}
 L(d)\,x_{\mathrm{phys}}(1+y_1)^{2d}
 <\frac{5875}{1296}\Bigl(1+\frac1{2d-1}\Bigr)^{2d}
 \le\frac{5875}{1296}\cdot\frac{256}{81}=\frac{94\,000}{6561}\approx14.3271<15 .
\end{equation}
Substituting \eqref{eq:coeff}, \eqref{eq:qsharp} and \eqref{eq:prefactor} into
\eqref{eq:geometric} yields, for all distinct $x,y\in\Lambda$,
\begin{align}
 \bigl|\langle\Dop_x^\dagger\Dop_y\rangle\bigr|
 &\le L(d)\,x_{\mathrm{phys}}(1+y_1)^{2d}\sum_{k\ge d_{\mathrm{dist}}}
 \Bigl(\frac{x_{\mathrm{phys}}}{x_1}\Bigr)^{k}\notag\\
 &<15\sum_{k\ge d_{\mathrm{dist}}}q_{\mathrm{sharp}}(d)^{k}
 =\frac{15}{1-q_{\mathrm{sharp}}(d)}\,q_{\mathrm{sharp}}(d)^{d_{\mathrm{dist}}}\notag\\
 &\le\frac{15}{1-q(d)}\,q_{\mathrm{sharp}}(d)^{d_{\mathrm{dist}}}
 =A(d)\,q_{\mathrm{sharp}}(d)^{d_{\mathrm{dist}}}\le A(d)\,q(d)^{d_{\mathrm{dist}}},
 \label{eq:decay}
\end{align}
with $A(d)=\frac{15}{1-q(d)}=\frac{360(d+1)}{14d+29}$. This proves
\eqref{eq:mainbound}.

\subsection{Uniform row summability}

On $\mathbb Z^d$ the distance $|z|_1=\sum_j|z_j|$ factors across coordinates, so
\begin{equation}\label{eq:rowsumderivation}
 \sum_{y\in\Lambda}\bigl|\langle\Dop_x^\dagger\Dop_y\rangle\bigr|
 \le1+A(d)\sum_{z\in\mathbb Z^d\setminus\{0\}}q(d)^{|z|_1}
 =1+A(d)\Bigl[\Bigl(\frac{1+q(d)}{1-q(d)}\Bigr)^{d}-1\Bigr]<\infty,
\end{equation}
the term $1$ coming from $y=x$, where the two insertions coincide and
$\|\Dop_x^\dagger\Dop_x\|_{\mathrm{op}}=1$. This proves \eqref{eq:rowsum}. The factor
$\bigl(\frac{1+q}{1-q}\bigr)^d$ is the only place where the sum over $\mathbb Z^d$
rather than a fixed finite window is used, and it is finite for every $d$ because
$q(d)<\frac5{12}<\frac12$; at $d=2$ it equals
$\bigl(\frac{1+5/24}{1-5/24}\bigr)^2=\frac{841}{361}$.

\section{Canonical thermodynamic limits}\label{sec:canonical}

Everything so far is grand canonical. The transfer to canonical states uses the
pair-source construction of Section~\ref{sec:kp} and the pressure of the \mbox{source-perturbed} ensemble, and it is a statement about thermodynamic limits, not about
uniform finite-volume canonical estimates.

Partition the torus into open hypercubes of side $\ell$ and keep only the bonds and
source pairs with both ends in the same block. This removes at most $dV/\ell$
boundary bonds, of unit operator norm, and at most $dV/\ell$ boundary source pairs,
each of norm at most $2\beta L|\xi|$, so their total logarithmic contribution per
site vanishes as $\ell\to\infty$ after $M\to\infty$. Because the local occupation is
bounded by two, the number-allocation entropy of distributing electrons among
$\ell^d$-site blocks is at most $\log(2\ell^d+1)/\ell^d$, or
$2\log(\ell^d+1)/\ell^d$ for separately fixed spin numbers, and both tend to zero.
The block decomposition therefore produces a concave canonical free-energy density
$f(\rho)$ and its Legendre dual $p(\mu)=\sup_{\rho\in[0,2]}[f(\rho)+\beta\mu\rho]$,
with $f$ concave and finite on $(0,2)$.

At every interior density $\rho\in(0,2)$ a finite supporting chemical potential
$\mu_*(\rho)$ exists, and by spin-exchange symmetry the balanced spin sectors
$(\rho/2,\rho/2)$ with $\rho/2\in(0,1)$ are included. Since the
Koteck\'y-Preiss convergence \eqref{eq:budgetsrc} on the source disk
\eqref{eq:disk} and the resulting bound \eqref{eq:decay} are uniform in
$\mu\in\mathbb R$ and in $M\ge4$, the secant-to-derivative comparison may be made
uniformly: on the one hand the grand-canonical correlation at $\mu$ is bounded by
\eqref{eq:decay}, on the other hand it is the derivative of the pressure with
respect to the source field, and the quadratic Taylor remainder of the \mbox{source-perturbed} pressure is bounded by $4\xi^2/\xi_0^2$ for $|\xi|\le\xi_0/2$, with
$\xi_0=\tau/(16L)$, because the expansion converges on the disk $\xi_0$ and Cauchy's
formula applies there. Choosing
the source amplitude of order $\ell^{-(d+1)}$ and letting $\ell\to\infty$ transfers
\eqref{eq:decay} and \eqref{eq:rowsum} with identical constants to every
translation-invariant canonical limit and to balanced separately fixed spin sectors.
This proves part 2 of Theorem~\ref{thm:main}.

The bounds are thus a statement about clustering of a two-point function, uniform in
$\mu$ and in the volume, and they pass to canonical limits. They say nothing about
finite-volume canonical states, about the existence of a phase transition, or about
which of several thermodynamic limits is realised. Part 3 of
Theorem~\ref{thm:main} follows because Definition~\ref{def:tpair} asks precisely for
the uniform decay and summability that \eqref{eq:decay} and \eqref{eq:rowsum}
supply, uniformly in $\Omega$.

\section{The baseline routes and the cubic obstruction}\label{sec:baselineproofs}

\subsection{Proof of Proposition~\ref{prop:a0}}

Take the connected bond-set grouping. For a connected bond set $E$ with $k=|E|$, every
vertex of $V(E)$ has at most $\Dphys(d)=2d$ incident bonds of $E$, and
$|V(E)|\le k+1$ because a connected graph on $k$ edges spans at most $k+1$ vertices.
The edge-local form \eqref{eq:localform} with $\alpha=U/(8d)$ gives, on the span of
$E$,
\[
 g\sum_{e\in E}B_e\le\frac U4\Nodd(V(E))+\frac{8\Dphys(d)\,k\,g^2}{U}I,
\]
with the handshake \eqref{eq:handshake} now applied inside $V(E)$. The atomic ratio
and the deletion factor contribute at most $(\frac54e)^{|V(E)|}\le(\frac54e)^{k+1}<\frac{15}{4}(\frac{15}{4})^k$,
using $e<\frac{1175}{432}<3$ and hence $\frac54e<\frac{15}{4}$. With the doubled-edge
Euler-tour count $a_k^{\mathrm{bond}}(G)\le\Daug(d)^{2k}$ for connected bond sets of
size $k$ in a graph of maximum degree $\Daug(d)$ \cite{rens2d}, the physical and
source step weights become
\begin{equation}\label{eq:a0weights}
 q_{\mathrm{phys}}^{\mathrm{A0}}(d)=\frac{15\Dphys(d)^2}{4L_{\mathrm{ET}}(d)}e^{\delta_{\mathrm{A0}}(d)},
 \qquad
 q_{\mathrm{src}}^{\mathrm{A0}}(d)=\frac{15\Daug(d)^2}{4L_{\mathrm{ET}}(d)}e^{\max(\delta_{\mathrm{A0}}(d),1/8)},
\end{equation}
with the $k=0$ site prefactor $\frac54e<\frac{15}{4}$. With
$L_{\mathrm{ET}}(d)=24\Daug(d)^2=96(d+1)^2$ and
$C_{\mathrm{ET}}(d)=64\Dphys(d)L_{\mathrm{ET}}(d)^2=128\,d\,L_{\mathrm{ET}}(d)^2$,
\begin{equation}\label{eq:a0delta}
 \delta_{\mathrm{A0}}(d)=\frac{8\Dphys(d)L_{\mathrm{ET}}(d)^2\beta t^2}{U}
 \le\frac{8\Dphys(d)}{64\Dphys(d)}=\frac18<\frac14,
 \qquad e^{\max(\delta_{\mathrm{A0}},1/8)}=e^{1/8}<e^{1/4}<\frac43,
\end{equation}
so
\begin{equation}\label{eq:a0decay}
 q_{\mathrm{src}}^{\mathrm{A0}}<\frac{15\Daug^2}{4L_{\mathrm{ET}}}\cdot\frac43
 =\frac{5\Daug^2}{L_{\mathrm{ET}}}=\frac5{24},
 \qquad
 \frac{15}{4}\cdot\frac{q_{\mathrm{src}}^{\mathrm{A0}}}{1-q_{\mathrm{src}}^{\mathrm{A0}}}
 <\frac{15}{4}\cdot\frac{5/24}{19/24}=\frac{75}{76}<1,
\end{equation}
and $q_{\mathrm{phys}}^{\mathrm{A0}}<\frac{5\Dphys^2}{L_{\mathrm{ET}}}=\frac{5d^2}{24(d+1)^2}=q_{\mathrm{ET}}(d)<\frac5{24}$.
The prefactor $\frac{15/4}{1-q_{\mathrm{ET}}(d)}$ is bounded above by
$\frac{15/4}{1-5/24}=\frac{90}{19}<\frac{76}{5}$; Proposition~\ref{prop:a0} quotes the
coarser value $\frac{12}{1-5/24}=\frac{288}{19}<\frac{76}{5}$, which dominates this bound.
This proves Proposition~\ref{prop:a0}.

The comparison with Section~\ref{sec:kp} is the point. Both routes use the same
edge-local estimate and the same odd-gap share, and the only difference is the count
of connected supports: $24\Daug^2$ against $36\Daug$ for the usage scale, and the
square of the scale cost against a single power. The Euler-tour route needs
$L\propto\Daug^2$ because it pays per bond, and its ceiling therefore carries
$\Daug^4\propto d^4$ times a further $d$ from $\Dphys=2d$ in \eqref{eq:a0weights},
which is $\Theta(d^5)$.

\subsection{Proof of Proposition~\ref{prop:b1}}

Use the conditional form \eqref{eq:cs} and \eqref{eq:b0} on a $k$-edge connected
bond set. Under $t/U\le\varepsilon_0(d)$, the form bound holds, and under
$\beta t^2/U\le1/C_{\mathrm{CF}}(d)$ the quadratic residual is
\begin{equation}\label{eq:b1residual}
 \frac{16\Dphys(d)\,k\,L_{\mathrm{CF}}(d)^2\,\beta t^2}{U}\le\frac k4,
\end{equation}
since $C_{\mathrm{CF}}(d)=64\Dphys(d)L_{\mathrm{CF}}(d)^2$. Together with the coarse
atomic and deletion constants of that route, the step weight at degree
$D\in\{\Dphys,\Daug\}$ is bounded by
\begin{equation}\label{eq:b1weight}
 \frac{D^2e^{9/4}}{L_{\mathrm{CF}}(d)}<\frac{27D^2}{L_{\mathrm{CF}}(d)},
\end{equation}
because $e<3$ and hence $e^{9/4}<27$. With
$L_{\mathrm{CF}}(d)=2^{10}\Daug(d)^2=4096(d+1)^2$,
\begin{equation}\label{eq:b1numbers}
 \varepsilon_0(d)=\frac{1}{16\Dphys L_{\mathrm{CF}}}=\frac{1}{2^{17}d(d+1)^2},\quad
 C_{\mathrm{CF}}(d)=64\Dphys L_{\mathrm{CF}}^2=2^{31}d(d+1)^4,
\end{equation}
the source activity is $\frac{27\Daug^2}{L_{\mathrm{CF}}}=\frac{27}{1024}$, the
Koteck\'y-Preiss budget is
\begin{equation}\label{eq:b1kp}
 e^2\,\frac{27/1024}{1-27/1024}<\frac{243}{997}<1
\end{equation}
since $e^2<9$, and the physical rate is
$q_{\mathrm{CF}}(d)=\frac{27\Dphys^2}{L_{\mathrm{CF}}}=\frac{27d^2}{1024(d+1)^2}$ with
prefactor $\frac{9}{1-q_{\mathrm{CF}}(d)}\le\frac{9216}{997}$. This proves
Proposition~\ref{prop:b1}.

\subsection{Proof of Proposition~\ref{prop:b2}}

Keep the conditional form, now on a connected vertex support $S$ with $v=|S|\ge2$.
Since $|E_{\mathrm{phys}}(S)|\le dv$, \eqref{eq:cs} gives
\[
 g\sum_{e\in E_{\mathrm{phys}}(S)}B_e\le\frac U4\Nodd(S)+\frac{32d^2g^2}{U}v\,I
\]
under \eqref{eq:smallcutoff}. With $L_{\mathrm{VA}}(d)=120\Daug(d)=240(d+1)$ and
$C_{\mathrm{VA}}(d)=\frac{512}{3}d^2L_{\mathrm{VA}}(d)^2$, the physical penalty per
site is $\delta_{\mathrm{B2}}(d)=\frac{32d^2L_{\mathrm{VA}}^2\beta t^2}{U}\le\frac{3}{16}$.
Choose the source disk $|\xi|\le\tau/(32L_{\mathrm{VA}})$, on which
$2\beta L_{\mathrm{VA}}|\xi|\le\frac1{16}$ and hence
$\delta_{\mathrm{tot}}^{\mathrm{B2}}\le\frac3{16}+\frac1{16}=\frac14$, giving
\begin{equation}\label{eq:b2K}
 L_{\mathrm{VA}}(d)\,x_{\mathrm{B2}}(d)<K_{\mathrm{B2}}:=\frac43\Bigl(\frac{1175}{432}\Bigr)^2
 =\frac{1\,380\,625}{139\,968}<\frac{1200}{121}<10 .
\end{equation}
Since $(11/10)^{11}=\frac{285\,311\,670\,611}{10^{11}}>\frac{1175}{432}>e$, evaluating
Lemma~\ref{lem:gf} at $y_2=1/(11\Daug(d))$ gives $(1+y_2)^{\Daug(d)}-1<\frac1{10}$ and
$x_{\mathrm{B2}}(d)<\frac{10}{121\Daug(d)}<x_0(y_2)$, so the Koteck\'y-Preiss budget is
\begin{equation}\label{eq:b2kp}
 \sum_{S\ni x_0,\,|S|\ge2}|w_{\mathrm{B2}}(S)|\,e^{|S|}
 <\frac{1\,380\,625}{139\,968}\cdot\frac1{10}=\frac{276\,125}{279\,936}\approx0.9864<1
\end{equation}
and the rate obeys
\begin{equation}\label{eq:b2q}
 q_{\mathrm{VA}}(d)<\frac{1\,380\,625}{139\,968}\cdot\frac{1175(2d-1)}{432\cdot240(d+1)}
 <\frac{1175(2d-1)}{10\,368(d+1)}<\frac{2d-1}{8(d+1)}<\frac14<\frac5{11} .
\end{equation}
This proves Proposition~\ref{prop:b2}. The two improvements act differently. Vertex
animals buy exactly one power of $d$ asymptotically, since
$C_{\mathrm{VA}}(d)=9\,830\,400\,d^2(d+1)^2$ against
$C_{\mathrm{ET}}(d)=1\,179\,648\,d(d+1)^4$, and they cost nothing when the conditional
form \eqref{eq:cs} is already available; note that $C_{\mathrm{VA}}(d)$ is in fact the
smaller of the two for $d\ge7$, where the two powers of $\Daug$ in the Euler-tour
count are no longer outweighed by the $\Dphys\propto d$ paid there. The edge-local
squares do something else: they remove the conditional cutoff entirely and, at $d=3$,
take the ceiling from $C_{\mathrm{VA}}(3)=1\,415\,577\,600$ to $C(3)=63\,700\,992$, a
further factor of $200/9$; against the naive cubic Euler-tour extension the total
gain at $d=3$ is the factor $128/9$ from $905\,969\,664$.

\subsection{Proof of Proposition~\ref{prop:c}}\label{sec:propc}

At $d=3$ the degrees are $\Dphys=6$ and $\Daug=8$. In the canonical formulation of
the Euler-tour route at side length parameter $L_7>0$, the trace penalty is quadratic
in the scale,
\begin{equation}\label{eq:d6}
 \delta_6=\frac{48L_7^2\beta t^2}{U},
\end{equation}
which is the value $3/8$ when $L_7=1024$ and $\beta t^2/U=2^{-27}$. The two step
weights of that route are rational multiples of $e^{\delta_6}$ with the coefficients
$135$ and $240$ at the physical and source degrees,
\begin{equation}\label{eq:qroutes}
 q_{\mathrm{phys}}=\frac{135}{L_7}e^{\delta_6},\qquad
 q_{\mathrm{src}}=\frac{240}{L_7}e^{\delta_6},
\end{equation}
and the canonical Koteck\'y-Preiss condition there reads
\begin{equation}\label{eq:canKP}
 \frac{15}{4}\,\frac{q_{\mathrm{src}}}{1-q_{\mathrm{src}}}\le1
 \quad\Longleftrightarrow\quad q_{\mathrm{src}}\le\frac{4}{19},
\end{equation}
since $15q\le4-4q$ rearranges to $19q\le4$. Now fix $C=2^{27}$ with
$\beta t^2/U=1/C$ and put $u=48L_7^2/C>0$, so that $\delta_6=u$. Then
\begin{equation}\label{eq:qsq}
 q_{\mathrm{src}}(L_7)^2=\frac{240^2}{L_7^2}e^{2u}
 =\frac{240^2\cdot48}{C}\cdot\frac{e^{2u}}{u}
 =\frac{675}{32\,768}\cdot\frac{e^{2u}}{u}.
\end{equation}
For every $u>0$ the quadratic Taylor polynomial of $e^{2u}$ is exact enough:
$1+2u+2u^2=4u+\frac12\bigl((2u-1)^2+1\bigr)>4u$, hence $e^{2u}>4u$ and
\begin{equation}\label{eq:lowerq}
 q_{\mathrm{src}}(L_7)^2>\frac{675}{32\,768}\cdot4=\frac{675}{8192}
 \approx0.08239>\frac{16}{361}=\Bigl(\frac4{19}\Bigr)^2\approx0.04432 .
\end{equation}
At $C=2^{27}$ this holds for every $L_7>0$, so no scale rescues the constant
$2^{27}$ in that route in dimension three, and the routes at $L_7=1024$ and
$L_7=720$, the latter being the retuned value tried in the earlier cubic extension,
both fail.

The failure is quantitative, not marginal. Minimizing $q_{\mathrm{src}}$ over $L_7$
at fixed $C$ is a one-variable problem: the logarithmic derivative $96L_7/C-1/L_7$
vanishes at $L_7^2=C/96$, where $\delta_6=\frac12$ and
\begin{equation}\label{eq:qmin}
 \min_{L_7>0}q_{\mathrm{src}}=240\sqrt{\frac{96e}{C}} .
\end{equation}
The condition $q_{\mathrm{src}}\le4/19$ then requires
\begin{equation}\label{eq:cmin}
 C\ge96\,e\cdot1140^2=124\,761\,600\,e\approx3.3914\times10^{8}=2^{28.337}>2^{28.33},
\end{equation}
whereas the two-dimensional value is $2^{27}$: the least admissible ceiling of that
route in three dimensions is $2^{1.337}$ times, that is about two and a half times,
the value it was asserted to keep.
The reason is visible in \eqref{eq:d6} and \eqref{eq:qroutes}: at $d=3$ the trace
penalty is $\frac{48L_7^2\beta t^2}{U}$ and the source weight carries
$240/L_7$, so the product is minimized when $\delta_6=\frac12$, and the value at that
optimum grows like $\sqrt{96e/C}$, which is far above $4/19$ at $C=2^{27}$. This
proves Proposition~\ref{prop:c}.

\section{Discussion}\label{sec:discussion}

\subsection{\texorpdfstring{Large-$d$ asymptotics}{Large-d asymptotics}}

Table~\ref{tab:asymptotics} collects the four rigorous routes. Replacing Euler-tour
counting by Bethe-tree vertex-animal counting improves the ceiling from
$\Theta(d^5)$ to $\Theta(d^4)$ and the conditional threshold from $\Theta(d^{-3})$ to
$\Theta(d^{-2})$; the edge-local squares then remove the threshold altogether
($\varepsilon_0=+\infty$).

\begin{table}[t]
\centering
\small
\setlength{\tabcolsep}{6pt}
\begin{tabular}{@{}lcccc@{}}
\hline
Route & $\varepsilon_0(d)$ as $d\to\infty$ & $C(d)$ as $d\to\infty$ & $q(d)$ & $A(d)$\\
\hline
Theorem~\ref{thm:main}  & $+\infty$ & $442\,368\,d^4(1+O(d^{-1}))$ & $\nearrow\frac5{12}=0.4167$ & $\nearrow\frac{180}7=25.71$\\
Proposition~\ref{prop:a0} & $+\infty$ & $1\,179\,648\,d^5(1+O(d^{-1}))$ & $\nearrow\frac5{24}=0.2083$ & $\frac{76}5=15.2$\\
Proposition~\ref{prop:b1} & $\frac1{2^{17}d^3}(1+O(d^{-1}))$ & $2^{31}d^5(1+O(d^{-1}))$ & $\nearrow\frac{27}{1024}=0.0264$ & $\le\frac{9216}{997}=9.24$\\
Proposition~\ref{prop:b2} & $\frac1{7680d^2}(1+O(d^{-1}))$ & $9\,830\,400\,d^4(1+O(d^{-1}))$ & $<\frac5{11}=0.4545$ & $O(1)$\\
\hline
\end{tabular}
\caption{Asymptotic behaviour of the four rigorous routes. The decay rate of
Theorem~\ref{thm:main} is stated for $q(d)=5(2d-1)/(24(d+1))$; the slightly better
$q_{\mathrm{sharp}}(d)$ is not shown.}
\label{tab:asymptotics}
\end{table}

Where do the four powers of $d$ in $C(d)$ come from? Take them one at a time.

\emph{The first factor of $d$} comes from the handshake in the odd-occupation
penalty. A single odd site can be shared by up to $\Dphys(d)=2d$ bonds, so the
per-bond parameter in \eqref{eq:alpha} must be $\alpha=U/(8d)$, and the residual
$\frac{2g^2}{\alpha}$ is multiplied by $2d$.

\emph{The second factor of $d$} comes from the maximum induced bond density
$|E_{\mathrm{phys}}(S)|\le d|S|$ in \eqref{eq:unconditional}, which turns
$\frac{16dg^2}{U}$ per bond into $\frac{16d^2g^2}{U}$ per site and hence
$\delta_{\mathrm{phys}}=\frac{16d^2L^2\beta t^2}{U}$.

\emph{The remaining factor of $d^2$} comes from the polymer entropy. The growth
constant of vertex animals on the source-extended graph is $\Daug(d)=2d+2$, so
suppressing the animal series requires a first-appearance factor
$L(d)=36\Daug(d)=\Theta(d)$, and since $g=L(d)t$ enters quadratically in
$\frac{2g^2}{\alpha}$, the scale contributes $L(d)^2=\Theta(d^2)$.

The product is $d\cdot d\cdot d^2=d^4$. The Euler-tour route has $L_{\mathrm{ET}}\propto\Daug^2$
and therefore $L_{\mathrm{ET}}^2\propto d^4$, whence $C_{\mathrm{ET}}\propto d\cdot d^4=d^5$.
The one power saved is precisely the difference between paying per bond and paying
per connected vertex support: a bond set of size $k$ spanning at least $k/2$ vertices
is counted once by the animal route, not $2^{|E(S)|}$ times.

\subsection{Heuristic comparison with the superexchange scale}

The following comparison is heuristic. It is not part of the theorem, no step of the
proof uses it, and it does not identify a transition or a material.

In the anti-adiabatic strong-coupling limit, second-order perturbation theory on the
doubly-occupied and empty pseudospin subspace assigns each nearest-neighbour bond the
superexchange coupling $J\sim4t^2/U$ \cite{langfirsov}. On a lattice of coordination
$z=\Dphys(d)=2d$ the Weiss scale of that coupling is
\begin{equation}\label{eq:heuristic}
 k_BT_{\mathrm{heur}}(d)\sim zJ=2d\cdot\frac{4t^2}{U}=8d\,\frac{t^2}{U}=\Theta(d)\frac{t^2}{U},
\end{equation}
or half of that in the standard Heisenberg normalisation of the exchange, which does
not change the order of growth. The rigorous ceiling grows like
$C(d)\approx442\,368\,d^4$, so
\begin{equation}\label{eq:ratio}
 \frac{C(d)}{8d}=55\,296\,d(d+1)^2\sim55\,296\,d^3 .
\end{equation}
What is shared with the heuristic scale is the dependence on the electron-phonon
coupling: both are proportional to $t^2/U$ and neither involves $\Omega$. What is not
shared is the growth with $d$, and the three extra powers are the three factors
traced above. They are the price of a sign-free majorant: any route that counts
connected supports on a degree-$2d$ graph by a hard-core cluster expansion pays
$L(d)=\Theta(d)$ per site, and no local rearrangement of the estimates removes that
factor, because the animal series genuinely grows with rate $\Theta(d)$ per site.
That a $\Theta(d^4)$ ceiling is not the true scale at large $d$ is therefore an
expectation, not a consequence of the present argument.

\subsection{Limitations}

The bound is explicit but not optimal, and in two respects it is deliberately loose.
The rate quoted in Theorem~\ref{thm:main} is $q(d)\to\frac5{12}$, while the sharper
value obtained in the proof obeys $q_{\mathrm{sharp}}(d)\to 6\,903\,125/20\,155\,392\approx0.3425$; the simpler
expression is kept because it gives the closed-form prefactor
$A(d)=360(d+1)/(14d+29)$. The constant $C(d)$ is not expected to be sharp for any
$d$, and the comparison $C(3)<2^{26}<2^{27}$ measures better bookkeeping rather than
a better physical scale.

The theorem bounds a clustering temperature of one two-point function, uniformly in
the frequency and in the chemical potential, in both ensembles, via a canonical
transfer through a source perturbation. It does not exclude order at temperatures
below the bound, it does not identify the order parameter of any phase, and it makes
no claim about a critical temperature or about a superconducting state. Physically,
the model is a caricature: it has no direct electron-electron interaction, no phonon
coupling beyond the on-site term, no hopping beyond nearest neighbours, and no disorder.

Three technical restrictions should be stated plainly. The oscillator problem is
treated at every $\Omega>0$ and no statement is made at $\Omega=0$, although the
bounds are uniform as $\Omega\downarrow0$ through positive values; the released
adiabatic bounds \cite{rensadiab,slow}, whose threshold vanishes as $\Omega\downarrow0$,
are sharper in that limit and are not improved here. The torus side
satisfies $M\ge4$, the condition under which the short-distance cycle structure of
$\mathbb Z^d$ used in Lemma~\ref{lem:animals}(3) is the square-lattice one. And the
canonical statement is a statement about thermodynamic limits at interior density,
compatible with balanced separately fixed spin sectors, and not about finite-volume
canonical states.

\section{Supporting material}\label{sec:supplement}

The preprint is archived at DOI \DOI{\PAPERDOI} and the supplement is archived at
DOI \DOI{\SUPPLEMENTDOI}.

The supplement is a single program of about seven hundred lines that recomputes, in
exact rational and symbolic arithmetic, every finite witness quoted above; it needs no
data file. Run without arguments it prints seventeen certified statements and exits
zero only if all of them hold, covering the rational enclosures that fix $e$ and its
powers, the single-bond and multi-bond form checks of Section~\ref{sec:forms}, the
animal counts of Lemma~\ref{lem:animals}, the constants and budgets of
Sections~\ref{sec:kp}, \ref{sec:decay}, \ref{sec:canonical}
and~\ref{sec:baselineproofs}, and the limits quoted in Section~\ref{sec:discussion}.
Run with the flag \verb|--mutants| it repeats the same checks with one guarded claim
corrupted at a time, fifteen times, and exits zero only if every corrupted claim is
rejected.

The certified run and the mutant run were carried out on a clean copy of the
supplement with uv 0.9.5, Python 3.13.9, sympy 1.14.0 and mpmath 1.3.0, and both
exited zero. The manifest in the supplement lists the SHA-256 digests of its two files;
the certificate program has digest
\texttt{c7360a785071fbb7afbd94723fae4866dbc530b4598c8f744c7e5a93187d030b}.

What the program certifies is arithmetic, not analysis. Its output lines assert exact
rational or symbolic identities and inequalities; they do not replace the estimates of Sections~\ref{sec:forms} to~\ref{sec:canonical},
which are ordinary inequalities for
operators and for real numbers with a written derivation. Where a displayed constant
can be checked by hand, as with $C(2)=15\,925\,248$ and $C(3)=63\,700\,992$, the
program agrees with the text.

\section*{AI assistance}

The analytical results, the design of the estimates, and the final text are the
author's responsibility. Independent proof reviews of the argument and of the
certificate program were carried out by separate AI systems, whose comments led to
several corrections recorded in the present version. The manuscript, its LaTeX
source, and the certificate program were drafted and edited with AI assistance under
the author's direction and review.
\addcontentsline{toc}{section}{AI assistance}

This preprint has not been peer reviewed.

\appendix

\section{Notation}\label{app:notation}

\begin{tabular}{@{}ll@{}}
$\Lambda$ & hypercubic torus $(\mathbb Z/M\mathbb Z)^d$, volume $V=M^d$\\
$\Dphys(d)$ & physical coordination number, $2d$\\
$\Daug(d)$ & source-extended degree, $2d+2$\\
$B_e$, $\Nodd(X)$, $\deg$ & bond transfer operator, odd-occupation count, degree\\
$g$ & bond-usage scaled hopping, $g=L(d)\,t$\\
$\delta_{\mathrm{phys}}(d)$ & physical per-site penalty, $\frac{16d^2L(d)^2\beta t^2}{U}\le\frac{3}{16}$\\
$\delta_{\mathrm{src}}(d)$ & source penalty, $\le\frac{1}{8}$\\
$L(d)$, $C(d)$, $q(d)$, $A(d)$ & scale, ceiling, decay rate, prefactor of Theorem~\ref{thm:main}\\
$x_{\mathrm{eff}}$, $x_0$, $x_1$ & uniform majorant of the activities, its Bethe limit, decay ratio\\
$T_{\mathrm{pair}}$ & pair-correlation temperature of Definition~\ref{def:tpair}\\
\end{tabular}

\bibliographystyle{plain}
\bibliography{refs}

\end{document}
